\chapter{The Neutron as a Composite Vortex and the Mechanics of Decay}


%----------------------------------------------------------------------
\section{Introduction: Beyond the Fundamental}
%----------------------------------------------------------------------

In the Standard Model, the neutron is a fundamental baryon composed of quarks, whose decay is governed by the Weak Nuclear Force.  In SDT, the neutron is not fundamental.  It is a \textbf{composite, mechanically-bound system}: a temporary, high-energy union of a proton and an electron.  Its neutrality, mass, and finite lifetime are direct, calculable consequences of SDT forces.


%----------------------------------------------------------------------
\section{Structure: A Packed Hydrogen Vortex}
%----------------------------------------------------------------------

\subsection{The SDT Hypothesis}

A neutron is formed by the extreme compression of a proton and an electron, their opposite charges cancelling to create a neutral entity.

\subsection{The Geometric Arrangement}

This is not a simple orbital system.  The immense pressure required (stellar cores, supernova shockwaves) forces the electron vortex to be \textbf{geometrically nested} within the proton's more powerful vortex.

\subsection{The Stabilising Lock (The Antineutrino)}

The packed configuration is not inherently stable.  The slight geometric and k-factor mismatch between proton and electron would cause immediate disintegration.  A third component---a topological knot in the spation field---is required as a ``lock.''  This stabilising gear \emph{is} the \textbf{antineutrino} ($\bar{\nu}_e$).

\begin{equation}\label{eq:neutron}
  n = [p^+ + e^-]_{\text{packed}} + \bar{\nu}_{e,\text{lock}}
\end{equation}


%----------------------------------------------------------------------
\section{The Mass of the Neutron: Energy of Compression}
%----------------------------------------------------------------------

The neutron is \emph{more} massive than the sum of its parts:
\begin{align}
  m_n &\approx 1.6749 \times 10^{-27}\;\text{kg} \\
  m_p + m_e &\approx 1.6735 \times 10^{-27}\;\text{kg} \\
  \Delta m &\approx 1.4 \times 10^{-30}\;\text{kg}
\end{align}

This mass difference is the energy \emph{input} required to overcome electrostatic repulsion and compress the electron vortex.  Stored in the stressed, spring-loaded configuration, it manifests as mass:
\begin{equation}\label{eq:binding}
  E_{\text{binding}} = \Delta m \cdot c^2 \approx 1.26 \times 10^{-13}\;\text{J} \approx 0.782\;\text{MeV}
\end{equation}


%----------------------------------------------------------------------
\section{Decay: A Violent Unzipping}
%----------------------------------------------------------------------

The free neutron's decay is not gentle.  It is a catastrophic mechanical failure.

\subsection{The Mechanism}

The neutron is an unstable equilibrium.  Internal stress from the mismatched proton and electron vortices, combined with external spation pressure, eventually overcomes the antineutrino lock.

\subsection{The Unzipping}

The lock fails.  The immense stored binding energy (0.782~MeV) is released instantly, converted into kinetic energy of the expelled components.  The neutron \textbf{disintegrates}.

\subsection{Ejection Velocities}

The explosive release propels the lighter electron and antineutrino to relativistic speeds:
\begin{equation}\label{eq:v-electron}
  v_{e,\text{max}} \approx 0.92\,c
\end{equation}

This is a direct prediction of the violent, mechanical nature of the decay event.


%----------------------------------------------------------------------
\section{The Electromagnetic Signature of Decay}
%----------------------------------------------------------------------

The unzipping is a catastrophic reconfiguration of the local spation medium.

\textbf{Initial state:} A single, compact, neutral displacement vortex.

\textbf{Final state:} Two separate, charged vortices flying apart at relativistic speeds.

This violent geometric change must create a powerful percussive shockwave.

\textbf{Prediction:} Every neutron decay event must be accompanied by a burst of high-energy photons---\textbf{gamma rays}---with energy peaked around the binding energy:
\begin{equation}\label{eq:gamma}
  E_{\gamma,\text{peak}} \approx E_{\text{binding}} = 0.782\;\text{MeV}
\end{equation}

Standard Model beta decay does not require an accompanying gamma ray.  Detecting this signature would be a \textbf{smoking gun} for SDT.


%----------------------------------------------------------------------
\section{The Lifetime and the ``Weak Force''}
%----------------------------------------------------------------------

The neutron's lifetime ($\sim$879\,s) is a measure of the \textbf{stability of the antineutrino lock}---the mean time to failure for this geometric configuration under constant spation pressure.

The Weak Interaction of the Standard Model is, in SDT, not a fundamental force.  It is a measure of \textbf{vortex instability and reconfiguration mechanics}.  It is ``weak'' not because the forces involved are small (the ejection velocities prove they are immense), but because the composite neutron vortex is remarkably stable, making catastrophic failure improbable over any short period.


%----------------------------------------------------------------------
\section{Conclusion}
%----------------------------------------------------------------------

By modelling the neutron as a composite proton-electron vortex stabilised by a geometric lock, we have:
\begin{enumerate}
  \item Provided a \textbf{mechanical structure} for the neutron.
  \item Accounted for its mass as constituents plus \textbf{energy of compression} (0.782~MeV).
  \item Described decay as \textbf{violent, deterministic ``unzipping''}.
  \item Predicted \textbf{relativistic ejection velocity} $v_e \approx 0.92\,c$.
  \item Made a novel, falsifiable prediction: a \textbf{$\sim$0.782~MeV gamma ray} for every decay event.
  \item Re-framed the ``Weak Force'' as \textbf{vortex instability}.
\end{enumerate}
